\chapter{D-Xylose Absorption Test}

\section*{What Is This Test?}
The D-xylose absorption test evaluates the absorptive capacity of the proximal small intestine. D-xylose is a pentose sugar that is absorbed in the duodenum and jejunum by passive carrier-mediated transport, requires no enzymatic digestion, and is excreted unchanged in urine. The test distinguishes proximal small-intestinal malabsorption from pancreatic enzyme insufficiency; in malabsorption from intestinal mucosal disease, oral D-xylose absorption is impaired, whereas in pancreatic insufficiency absorption is preserved (because the sugar does not require pancreatic enzymes). The test is largely historical and often replaced by direct measures of mucosal disease (biopsy) and pancreatic function (fecal elastase), but is still occasionally used.

\section*{Alternative Names}
\begin{itemize}
  \item D-Xylose Test
  \item Xylose Absorption Test
  \item Xylose Tolerance Test
  \item D-Xylose Absorption Test
\end{itemize}
\section*{Principle}
After an overnight fast, the patient drinks a solution of 25 g D-xylose. Urine is collected for 5 hours, and blood samples are drawn at 1 and 2 hours. In healthy individuals, blood xylose rises $>$ 20 mg/dL at 1 or 2 hours, and $>$ 4 g of xylose is excreted in 5 hours. Impaired absorption lowers blood levels and urinary excretion. The 1-hour blood level correlates well and is preferred in children.

\section*{History}
The D-xylose test was developed by Benson et al.\ in the 1950s and became the standard for distinguishing sprue from pancreatic insufficiency in the 1960s. Modern tests (anti-tTG, anti-endomysial antibodies, small bowel biopsy, fecal elastase) have largely replaced routine use.

\section*{Why This Test Is Ordered}
\begin{itemize}
  \item Distinguish small-intestinal mucosal disease (celiac disease, tropical sprue) from pancreatic exocrine insufficiency (chronic pancreatitis, cystic fibrosis)
  \item Evaluate unexplained malabsorption or weight loss when screening antibody tests are negative
  \item Monitor response to gluten-free diet in treated celiac disease
\end{itemize}

\section*{Primary vs Secondary Test}
Secondary test, performed when first-line evaluations (celiac serology, stool elastase, fecal fat) are inconclusive.

\section*{How This Test Is Performed}
Patient fasts overnight. Empty bladder at start. Drink 25 g D-xylose in 250 mL water, followed by additional fluids to ensure adequate urine output. Blood drawn at 1 and 2 hours. Urine collected over 5 hours. Patient remains supine and avoids eating. Vomiting invalidates the test.

\section*{What Preparation Is Needed}
Overnight fast; no significant fluid restriction.

\section*{Contraindications and Precautions}
\begin{itemize}
  \item Contraindicated in renal failure (false low urine excretion despite normal absorption); in such cases, blood levels alone are interpreted. Vomiting after xylose ingestion invalidates test and may indicate rapid rise in osmotic load. Diabetes mellitus, pregnancy, ascites may affect interpretation. Elderly patients have lower normal excretion.
\end{itemize}
\section*{Risks}
Nausea, abdominal discomfort, diarrhea from the osmotic load. Rare: hypoglycemia in diabetics.

\section*{What Happens After the Test}
Interpret with renal function. Reduced absorption indicates proximal small intestinal disease; further workup includes celiac serology, duodenal biopsy, capsule endoscopy.

\section*{Understanding Results}

\subsection*{Below Normal}
\begin{itemize}
  \item \textbf{Decreased blood xylose or 5-h urine $<$ 4 g:} Proximal small-intestinal malabsorption
  \item \textbf{Causes:} Celiac disease, tropical sprue, Crohn's disease of proximal small bowel, Whipple disease, eosinophilic gastroenteritis, small bowel bacterial overgrowth, lymphoma
  \item \textbf{Pancreatic insufficiency:} Normal xylose absorption (distinguishing feature)
  \item \textbf{Renal dysfunction}: False low urine excretion (use blood level)
  \item \textbf{Elderly}: Lower normal values
\end{itemize}

\subsection*{Above Normal}
Normal; no clinical significance except excluding malabsorption.

\section*{Normal Results}
\begin{itemize}
  \item \textbf{1- or 2-hour serum D-xylose (adults):} $>$ 20 mg/dL
  \item \textbf{5-hour urine D-xylose excretion (adults):} $>$ 4 g of 25 g dose (16\%)
  \item \textbf{Children 1-hour serum:} $>$ 20 mg/dL (with modified 5 g dose in younger children)
  \item \textbf{Renal insufficiency:} Use blood level only; urine excretion unreliable
  \item \textbf{Elderly $>$ 65 years:} Threshold may be $>$ 3.5 g/5 h
\end{itemize}

\section*{Follow-up Tests}
\begin{itemize}
  \item \textbf{Tissue transglutaminase IgA and total IgA}: Celiac disease screening
  \item \textbf{Duodenal biopsy with villous morphology}: Strongyloides, celiac disease confirmation
  \item \textbf{Fecal elastase-1 or 72-hour fecal fat}: Pancreatic insufficiency
  \item \textbf{Serum albumin, calcium, INR, vitamin levels (B12, folate, D, K, iron panel)}: Nutritional consequence of malabsorption
  \item \textbf{Small bowel imaging, capsule endoscopy}: For Crohn's disease or tumors
  \item \textbf{Lactose hydrogen breath test}: Secondary lactase deficiency
\end{itemize}

\section*{References}
\begin{enumerate}
  \item Casellas F, et al. D-xylose absorption test in clinical practice. Gastroenterol Hepatol. 2013;36(7):487-492.
  \item Benson JA Jr, et al. The d-xylose absorption test in malabsorption syndromes. N Engl J Med. 1957;256(8):355-359.
  \item Hovelius B, et al. D-xylose absorption in health and disease. Scand J Gastroenterol. 1986;21(9):1047-1053.
  \item Schiller LR, et al. Diarrhea and malabsorption. In: Feldman M, et al., eds. Sleisenger and Fordtran's Gastrointestinal and Liver Disease. 11th ed. Elsevier; 2021.
\end{enumerate}

\clearpage
\section*{Regulatory Status and Authorization Date}
This chapter describes a test category, not a specific commercial product. FDA, CDSCO, EU IVDR, and MHRA approval or registration dates therefore cannot be assigned without the assay manufacturer, product name, instrument, and intended use. For laboratory-developed tests, an individual agency approval date may not exist. See the Preface for the interpretation of regulatory status.

\clearpage
